\section{本章小结}
\label{sec:adjacency-matrix-summary}

邻接矩阵的$k$次幂在位置$(u,v)$记录从$u$到$v$的$k$步行走条数。对角元之和（迹）给出闭行走总数。实对称性保证迹等于特征值$k$次幂之和，这就是谱矩。

典型图的谱可以直接构造。完全图利用常向量与零和向量，完全二部图利用两侧对称，环利用周期边界的单位根，路径利用零端边界的正弦模式。笛卡儿积图的特征值等于各因子特征值之和。正则图的常向量特征值等于度，补图的谱由此换算。

谱本身具有结构性质。连通图的最大特征值是单根，有严格正的Perron向量。它夹在平均度和最大度之间，连通正则图中两端同时取等。删边或删正权顶点不增大谱半径。二部图的谱关于零对称，非二部图没有这种对称。

完整谱与图结构的关系体现在两个方面。一方面，谱记录了图的宏观代数特征：低阶谱矩确定边数和三角形数，极端特征值限制直径、独立数和色数。另一方面，谱丢失了同构信息：$K_{1,4}$和$C_4\sqcup K_1$有相同谱却不同构，不同来源的闭行走在求和后互相补偿。

更粗的商谱进一步缩减信息。等价划分把图缩成商矩阵，颜色细化给出最粗的等价划分。商谱保留了块间连接密度的框架，丢失了块内的微观差异——颜色细化分不出的两张图，完整谱仍可以不同。

下一章把邻接作用换成沿边差分再加权汇总，得到拉普拉斯矩阵。常值函数的沿边差分全为零，因此零空间直接揭示连通分量——恰好弥补邻接谱在这方面的缺失。
